\documentclass[12pt,oneside]{book}

% ============================================================
% METADATOS
% ============================================================
\newcommand{\universidad}{Universidad de Buenos Aires (UBA)}
\newcommand{\facultad}{Facultad de Derecho}
\newcommand{\carrera}{}
\newcommand{\materia}{Derecho de la Integración}
\newcommand{\titulo}{Jerarquía normativa, derechos humanos y derecho derivado en el Mercosur}
\newcommand{\autor}{Isaac Edgar Camacho Ocampo}
\newcommand{\anio}{2026}

% ============================================================
% PAQUETES
% ============================================================
\usepackage[utf8]{inputenc}
\usepackage[T1]{fontenc}
\usepackage[spanish,es-nodecimaldot]{babel}

\usepackage[
    top=2.5cm,
    bottom=2.5cm,
    left=3cm,
    right=2.5cm,
    headheight=15pt
]{geometry}

\usepackage{amsmath}
\usepackage{amssymb}
\usepackage{mathtools}

\usepackage{graphicx}
\usepackage{float}
\usepackage{caption}
\usepackage{subcaption}

\usepackage{booktabs}
\usepackage{array}
\usepackage{longtable}
\usepackage{tabularx}

\usepackage{enumitem}

\usepackage{xcolor}
\usepackage[most]{tcolorbox}

\usepackage{fancyhdr}
\usepackage{ragged2e}

\usepackage[
    colorlinks=true,
    linkcolor=blue!50!black,
    citecolor=green!40!black,
    urlcolor=blue!60!black,
    pdftitle={\titulo},
    pdfauthor={\autor},
    pdfsubject={Apuntes universitarios},
    bookmarks=true
]{hyperref}

% ============================================================
% RUTAS DE IMÁGENES
% ============================================================
\graphicspath{
    {./img/}
    {./graficos/}
    {./figuras/}
}

% ============================================================
% CONFIGURACIÓN GENERAL
% ============================================================
\setcounter{tocdepth}{2}

\setlength{\parindent}{0pt}
\setlength{\parskip}{0.5em}

\setlist[itemize]{
    topsep=0.3em,
    itemsep=0.2em
}

\setlist[enumerate]{
    topsep=0.3em,
    itemsep=0.2em
}

\clubpenalty=10000
\widowpenalty=10000

% ============================================================
% COMANDOS MATEMÁTICOS
% ============================================================
\newcommand{\abs}[1]{\left|#1\right|}
\newcommand{\norm}[1]{\left\lVert#1\right\rVert}
\newcommand{\vect}[1]{\mathbf{#1}}
\newcommand{\deriv}[2]{\frac{d#1}{d#2}}
\newcommand{\pderiv}[2]{\frac{\partial#1}{\partial#2}}

% ============================================================
% ENTORNOS / CAJAS ACADÉMICAS
% ============================================================
\newtcolorbox{definition}[1]{
    enhanced,
    breakable,
    title={Definición: #1},
    fonttitle=\bfseries,
    colback=blue!3,
    colframe=blue!50!black,
    boxrule=0.8pt,
    arc=2mm
}

\newtcolorbox{important}{
    enhanced,
    breakable,
    title={Importante},
    fonttitle=\bfseries,
    colback=yellow!8,
    colframe=orange!70!black,
    boxrule=0.8pt,
    arc=2mm
}

\newtcolorbox{example}[1]{
    enhanced,
    breakable,
    title={Ejemplo: #1},
    fonttitle=\bfseries,
    colback=green!4,
    colframe=green!40!black,
    boxrule=0.8pt,
    arc=2mm
}

\newtcolorbox{formula}[1]{
    enhanced,
    breakable,
    title={Fórmula / Regla: #1},
    fonttitle=\bfseries,
    colback=purple!4,
    colframe=purple!50!black,
    boxrule=0.8pt,
    arc=2mm
}

\newtcolorbox{exercise}[1]{
    enhanced,
    breakable,
    title={Ejercicio: #1},
    fonttitle=\bfseries,
    colback=gray!5,
    colframe=gray!60!black,
    boxrule=0.8pt,
    arc=2mm
}

\newtcolorbox{note}{
    enhanced,
    breakable,
    title={Nota},
    fonttitle=\bfseries,
    colback=white,
    colframe=gray!50,
    boxrule=0.6pt,
    arc=2mm
}

% ============================================================
% ENCABEZADOS Y PIES
% ============================================================
\pagestyle{fancy}
\fancyhf{}

\fancyhead[L]{\nouppercase{\leftmark}}
\fancyhead[R]{\materia}

\fancyfoot[C]{\thepage}

\renewcommand{\headrulewidth}{0.4pt}
\renewcommand{\footrulewidth}{0pt}

\fancypagestyle{plain}{
    \fancyhf{}
    \fancyfoot[C]{\thepage}
    \renewcommand{\headrulewidth}{0pt}
}

% ============================================================
% DOCUMENTO
% ============================================================
\begin{document}

% ------------------------------------------------------------
% PORTADA
% ------------------------------------------------------------
\begin{titlepage}

\thispagestyle{empty}

\begin{center}

\vspace*{1cm}

{\large \textsc{\universidad}}

\vspace{0.2cm}

{\large \textsc{\facultad}}

\vspace{3cm}

{\Huge \textbf{\materia}}

\vspace{1cm}

{\LARGE \titulo}

\vspace{0.8cm}

\textit{Comprender la jerarquía de las normas es aprender a distinguir, en cada conflicto jurídico, cuál es la voz que finalmente decide.}

\vspace{1.2cm}

\rule{0.70\textwidth}{0.5pt}

\vspace{0.5cm}

{\large Apuntes de estudio}

\vspace{0.5cm}

\rule{0.70\textwidth}{0.5pt}

\vfill

{\large \autor}

\vspace{0.3cm}

{\large \anio}

\end{center}

\vfill

\begin{flushleft}
\textit{Este es un modesto aporte a los estudiantes de ``Derecho de la Integración'' no pretende sustituir el material oficial de la materia.}
\end{flushleft}

\end{titlepage}

% ------------------------------------------------------------
% ÍNDICE
% ------------------------------------------------------------
\tableofcontents

% ============================================================
% CONTENIDO
% ============================================================

\chapter[Jerarquía normativa y derecho derivado]{Jerarquía normativa, derechos humanos y derecho derivado en el Mercosur}

La clase reconstruye el modo en que se ordenan, dentro del sistema jurídico argentino, la Constitución Nacional, los tratados internacionales, el derecho derivado de los procesos de integración —como el Mercosur— y las leyes internas. Se recorre la evolución jurisprudencial de la Corte Suprema de Justicia de la Nación —desde el artículo 31 de la Constitución Nacional hasta la reforma de 1994— a través de dos precedentes centrales (``Ekmekdjian c. Sofovich'' y ``Cafés La Virginia''), y se analizan los principios propios del derecho de la integración y del derecho comunitario: especificidad, atribución de competencias exclusivas y vigencia simultánea.

Determinar qué norma prevalece cuando dos disposiciones se contradicen —un contrato internacional, una resolución aduanera, una sentencia que invoca un tratado de derechos humanos— resulta indispensable para el ejercicio profesional del derecho, la asesoría a empresas que operan en el marco del Mercosur y la argumentación ante tribunales que invocan tratados internacionales.

El desarrollo de la clase puede describirse como la reconstrucción progresiva de una pirámide normativa. En su base se sitúa el artículo 31 de la Constitución Nacional y la distinción entre monismo y dualismo (Bloque 1); sobre esa base se apoya un caso concreto, ``Ekmekdjian c. Sofovich'', que puso a prueba esa pirámide mediante un conflicto real entre la libertad de expresión y el derecho a réplica (Bloque 2); ese conflicto impulsó el rediseño de la pirámide en la reforma constitucional de 1994 (Bloque 3); una vez reconstruida, se verifica su funcionamiento con otro caso, ``Cafés La Virginia'', y con el Mercosur (Bloque 4); y finalmente se exponen los principios generales que sostienen cualquier proceso de integración, sea el Mercosur o la Unión Europea (Bloques 5 y 6).

\section{Jerarquía normativa y sistemas de recepción del derecho internacional}

\subsection{Monismo y dualismo}

La distinción entre los sistemas dualista y monista de recepción del derecho internacional constituye el punto de partida necesario para analizar la jerarquía normativa. Discutir la jerarquía entre un tratado y una ley sólo tiene sentido dentro de un sistema monista, en el cual ambas categorías de normas pertenecen a un mismo ordenamiento y pueden compararse entre sí. En un sistema dualista, en cambio, el derecho internacional y el derecho interno se mueven en ámbitos distintos que no se cruzan, por lo que la pregunta por la jerarquía entre tratado y ley carece de sentido. La Argentina se ubica, a estos efectos, dentro de un sistema monista.

\subsection{El artículo 31 de la Constitución Nacional}

Más allá del artículo 75, inciso 22 —incorporado recién en la reforma de 1994—, la disposición constitucional que desde el origen mismo del texto de 1853 se refirió a la jerarquía normativa es el artículo 31.

\begin{definition}{Artículo 31 de la Constitución Nacional}
``Esta Constitución, las leyes de la Nación que en su consecuencia se dicten por el Congreso y los tratados con las potencias extranjeras son la ley suprema de la Nación; y las autoridades de cada provincia están obligadas a conformarse a ella [\ldots]''

\textit{Constitución Nacional Argentina, artículo 31.}
\end{definition}

A partir de esta disposición se planteó una extensa discusión doctrinaria: si la Constitución, las leyes dictadas por el Congreso y los tratados constituían, en conjunto, la ley suprema de la Nación sin distinción de orden interno, o si existía entre ellos una jerarquía determinada. La doctrina nunca puso en duda que la Constitución ocupara la cúspide de la pirámide normativa; el debate se centró, en cambio, en la relación entre la ley y los tratados.

\subsection{Evolución jurisprudencial: de la doctrina de guerra y paz al fallo ``Martín''}

Antes de la década de 1960, el criterio de la Corte Suprema de Justicia de la Nación distinguía según la Argentina se encontrara en épocas de paz o de guerra: en tiempos de paz se consideraba que el país se movía dentro de un sistema dualista, y sólo en momentos de guerra los tratados internacionales —en particular los referidos a la protección de las personas— adquirían relevancia jurídica; fuera de esos supuestos, únicamente se aplicaba la ley interna.

Ese criterio comenzó a modificarse a partir del fallo ``Martín c. Administración General de Puertos'' (Corte Suprema de Justicia de la Nación, 1963).

\begin{important}
La Corte Suprema sostuvo, con relación al artículo 31 de la Constitución Nacional, que la norma no distingue entre la ley y el tratado: llama a ambos ``ley suprema de la Nación''. El hecho de que el artículo mencione primero uno y después el otro no responde a una jerarquía, sino a una cuestión de redacción. Si la mención en primer lugar de ``las leyes'' implicara jerarquía, se estaría ante un sistema de monismo con supremacía del derecho interno absoluto, en el que las normas de fuente interna prevalecen siempre sobre el tratado. A partir de ``Martín'', la Corte entendió, en cambio, que la ley y el tratado tenían equivalencia jerárquica: se ubicaban en el mismo nivel de la pirámide normativa.
\end{important}

\subsection{Criterios para resolver conflictos normativos: ley especial y ley posterior}

Si el tratado y la ley tienen igual jerarquía, corresponde determinar cómo se resuelve un conflicto cuando ambos regulan de manera distinta una misma situación. La dogmática jurídica clásica interpreta al legislador como una entidad única y coherente, que no puede establecer soluciones contradictorias para un mismo supuesto de hecho. Cuando existen dos normas de igual jerarquía que se contradicen, la doctrina recurre a dos criterios de interpretación.

\begin{table}[H]
\centering
\caption{Criterios para resolver conflictos entre normas de igual jerarquía}
\begin{tabularx}{\textwidth}{@{}>{\raggedright\arraybackslash}p{0.28\textwidth}X@{}}
\toprule
\textbf{Criterio} & \textbf{Regla y efecto sobre la norma general} \\
\midrule
Ley especial prevalece sobre la general &
La norma que regula específicamente una materia se aplica por sobre la norma general que también podría abarcarla. No la deroga: la norma general sigue vigente para todos los demás supuestos no cubiertos por la ley especial. \\
\addlinespace
Ley posterior deroga la anterior &
Entre dos normas de igual generalidad, la dictada después reemplaza a la anterior. Sí la deroga: la norma anterior deja de tener vigencia en lo que resulte contradictorio. \\
\bottomrule
\end{tabularx}
\end{table}

\begin{example}{El contrato de compraventa internacional}
Si existe un contrato de compraventa internacional de mercaderías y, al mismo tiempo, rige una convención internacional de compraventa, esta última prevalece frente al capítulo de contratos del Código Civil, porque esa convención resulta especial para esa materia. La circunstancia de que una norma se llame ``tratado'' y la otra ``ley'' resulta indiferente a estos efectos, ya que ambas tienen igual jerarquía. La convención no deroga al Código Civil, que sigue rigiendo para la compraventa de mercadería no alcanzada por la convención, para la compraventa de transporte y para todos los demás supuestos.
\end{example}

En el período posterior al fallo ``Martín'', el criterio de la Corte fue de igualdad jerárquica entre tratado y ley, por lo que los conflictos entre ambos se resolvían con los mismos criterios de ley especial y ley posterior aplicables entre dos leyes.

\section{El caso ``Ekmekdjian c. Sofovich'' y los derechos humanos}

\subsection{Los hechos del caso}

Los hechos que dieron origen a este precedente ocurrieron en 1988, antes de la reforma constitucional de 1994, motivo por el cual a este caso se le sigue aplicando el artículo 31 de la Constitución Nacional.

\begin{example}{El caso ``Ekmekdjian c. Sofovich'' (Corte Suprema de Justicia de la Nación, 1992)}
En un programa de televisión conducido por Gerardo Sofovich, el escritor Dalmiro Sáenz realizó, en el marco de un debate sobre religión, manifestaciones que cuestionaban la virginidad de María y la condición de Jesús como hijo de Dios. Un profesor de derecho constitucional, Miguel Ángel Ekmekdjian, se sintió agraviado en su fe católica y solicitó a Sofovich, primero en forma privada y luego judicialmente, ejercer el derecho a réplica en el mismo programa.

Sofovich se negó, invocando su libertad de expresión: era su programa y decidía quiénes tenían lugar en él. No se trataba de una simple ofensa personal a la Virgen María, sino de un cuestionamiento a un principio básico de la fe católica —que Jesús es hijo de Dios, parte de la Santísima Trinidad—, lo cual explica que Ekmekdjian considerara vulnerado un derecho fundamental y no simplemente una susceptibilidad personal.
\end{example}

\subsection{Los derechos en pugna: libertad de expresión y derecho a réplica}

Sofovich invocaba la libertad de expresión; Ekmekdjian, el derecho a réplica.

\begin{definition}{Artículo 14 de la Constitución Nacional}
``Todos los habitantes de la Nación gozan de los siguientes derechos conforme a las leyes que reglamenten su ejercicio; a saber: [\ldots] de publicar sus ideas por la prensa sin censura previa [\ldots]''

\textit{Constitución Nacional Argentina, artículo 14.}
\end{definition}

En 1853 el único medio de comunicación existente era la prensa escrita, por lo que la fórmula ``por la prensa'' fue luego interpretada extensivamente para abarcar cualquier medio de comunicación masivo, incluidos la radio y la televisión. El problema radicaba en que el derecho a réplica no está mencionado expresamente en el artículo 14, sino que se fundamenta en el artículo 33.

\begin{definition}{Artículo 33 de la Constitución Nacional (derechos implícitos)}
``Las declaraciones, derechos y garantías que enumera la Constitución no serán entendidos como negación de otros derechos y garantías no enumerados; pero que nacen del principio de la soberanía del pueblo y de la forma republicana de gobierno.''

\textit{Constitución Nacional Argentina, artículo 33.}
\end{definition}

No todos los derechos están enumerados expresamente en el texto constitucional: existen derechos y garantías implícitos, que surgen del sistema republicano y democrático de gobierno aun sin estar enunciados, y en ese marco podría encuadrarse el derecho a réplica.

\paragraph{Reconocimiento, protección y promoción de derechos.}
Corresponde distinguir los tipos de obligaciones que puede asumir el Estado respecto de un derecho.

\begin{table}[H]
\centering
\caption{Tipos de obligaciones del Estado respecto de un derecho}
\begin{tabularx}{\textwidth}{@{}>{\raggedright\arraybackslash}p{0.2\textwidth}>{\raggedright\arraybackslash}p{0.35\textwidth}X@{}}
\toprule
\textbf{Tipo de obligación} & \textbf{Qué hace el Estado} & \textbf{Ejemplo} \\
\midrule
Reconocimiento & Declara que la persona ya es titular del derecho, sin agregar nada más. & Reconocer el derecho a la educación: ``eduque como quiera''. \\
\addlinespace
Protección & Además de reconocerlo, crea mecanismos para que nadie impida su ejercicio. & No discriminación dentro de las escuelas: garantizar el acceso al derecho. \\
\addlinespace
Promoción & Además de reconocerlo y protegerlo, el Estado facilita activamente su ejercicio. & Otorgar una beca de estudios para facilitar el acceso a la educación. \\
\bottomrule
\end{tabularx}
\end{table}

La Constitución reconoce la libertad de culto en general, pero además asume una obligación distinta y más intensa respecto del culto católico.

\begin{definition}{Artículo 2 de la Constitución Nacional}
``El Gobierno federal sostiene el culto católico apostólico romano.''

\textit{Constitución Nacional Argentina, artículo 2.}
\end{definition}

Se trata de dos tipos de obligaciones distintas: el Estado sostiene el culto católico apostólico romano y lo impulsa activamente —de allí que existan becas para escuelas católicas y no para escuelas de otras confesiones—, no porque se prohíban las demás, sino porque el Estado asumió una obligación de promoción exclusivamente respecto de la religión católica, que era la fe preponderante en la sociedad argentina de aquel momento.

\subsection{El Pacto de San José de Costa Rica y la Convención de Viena}

Descartado que el derecho a réplica surgiera expresamente del artículo 14, Ekmekdjian lo fundó además en una fuente convencional específica.

\begin{definition}{Convención Americana sobre Derechos Humanos --- Pacto de San José de Costa Rica (1969), artículo 14.1}
Reconoce el derecho de toda persona afectada por informaciones inexactas o agraviantes, emitidas en su perjuicio a través de medios de difusión, a efectuar por el mismo órgano su rectificación o respuesta, en las condiciones que establezca la ley.

\textit{Convención Americana sobre Derechos Humanos, artículo 14.}
\end{definition}

La Argentina ya había ratificado el Pacto de San José de Costa Rica al momento de los hechos; lo que no existía era una ley interna que reglamentara el modo de ejercer el derecho a réplica. Por ello se discutía si se trataba de un derecho operativo o meramente programático: el pacto contempla el derecho a réplica, pero lo somete a las pautas de su reglamentación, y el Estado argentino no la había dictado. Técnicamente, entonces, el actor no podría ejercerlo porque no existía un procedimiento establecido para hacerlo. Frente a esta dificultad, la Corte recurrió a otro tratado: la Convención de Viena sobre el Derecho de los Tratados, que establece las normas generales aplicables a cualquier tratado —cómo entra en vigor, cómo se denuncia, cómo se aplica— a falta de previsión expresa.

\begin{definition}{Convención de Viena sobre el Derecho de los Tratados (1969), artículo 27}
``Una parte no podrá invocar las disposiciones de su derecho interno como justificación del incumplimiento de un tratado.''

\textit{Convención de Viena sobre el Derecho de los Tratados, artículo 27 --- principio de ``pacta sunt servanda''.}
\end{definition}

Este es el principio de \textbf{pacta sunt servanda}: lo pactado es obligatorio. La Corte sostuvo, sobre esta base, que la Argentina estaba incumpliendo una obligación asumida, puesto que tampoco es posible invocar la propia falta de reglamentación para dejar de cumplir un tratado.

\subsection{La solución de la Corte y su trascendencia}

Ante la falta de reglamentación del derecho a réplica, la Corte Suprema resolvió, para este caso puntual, una reglamentación provisoria: reconoció a Ekmekdjian el derecho a réplica en las mismas condiciones del programa original —el mismo horario, el mismo tiempo de emisión—. En los hechos, Sofovich admitió la carta de descargo escrita por Ekmekdjian, mostrando en pantalla su texto durante unos minutos del programa, sin reabrir el debate, pero reconociendo el derecho a réplica sobre la base de una fuente convencional.

A partir de este fallo se sostiene que ``Ekmekdjian c. Sofovich'' invirtió el orden Constitución–Tratado–Ley que regía desde ``Martín'', en la medida en que la Corte recurrió a un tratado internacional (la Convención de Viena) para sostener que no podía dejarse sin efecto un derecho reconocido en otro tratado por la falta de reglamentación de una ley interna.

\begin{note}
Existían en este caso dos derechos en pugna: uno reconocido en el Pacto de San José de Costa Rica y otro en la propia Constitución Nacional. El artículo 33 establece que la declaración de derechos y garantías enumerados en la Constitución no puede invocarse como desconocimiento de otros derechos no numerados expresamente, pero que surgen de la forma republicana de gobierno y del estado de derecho —es decir, los derechos y garantías implícitos—. Bajo esta segunda lectura, el derecho a réplica podría fundarse directamente en el artículo 33 de la Constitución Nacional, sin necesidad de recurrir a la jerarquía de los tratados. En rigor, tratándose de un caso de derechos humanos, este precedente no bastaría por sí solo para sostener que todo tratado prevalece sobre toda ley, porque la solución también podía explicarse dentro del propio texto constitucional.
\end{note}

\section{La reforma constitucional de 1994 y la pirámide normativa argentina}

\subsection{El bloque de constitucionalidad}

La reforma constitucional de 1994 vino a resolver expresamente la discusión sobre la jerarquía entre tratado y ley. Lo primero que hizo fue incorporar los tratados de derechos humanos a la parte dogmática de la Constitución Nacional, mediante el artículo 75, inciso 22.

\begin{definition}{Artículo 75, inciso 22 de la Constitución Nacional (reforma de 1994)}
Enumera un conjunto de declaraciones y tratados de derechos humanos que ``en las condiciones de su vigencia, tienen jerarquía constitucional, no derogan artículo alguno de la primera parte de esta Constitución y deben entenderse complementarios de los derechos y garantías por ella reconocidos''. Los demás tratados de derechos humanos pueden alcanzar jerarquía constitucional con el voto de las dos terceras partes de la totalidad de los miembros de cada Cámara.

\textit{Constitución Nacional Argentina, artículo 75, inciso 22.}
\end{definition}

Este conjunto normativo constituye lo que un sector de la doctrina denomina la capa de constitucionalidad: el texto de la Constitución incorpora todos los tratados de derechos humanos allí enumerados, así como los que se ratifiquen posteriormente mediante ese procedimiento agravado. Una consecuencia práctica de esta incorporación es que, para denunciar uno de estos tratados, se exige un quórum agravado, equiparable al de una reforma constitucional, precisamente porque denunciarlo equivale a modificar la Constitución misma.

\subsection{La jerarquía superior de los tratados en general}

Después de los tratados de derechos humanos, el resto de los tratados internacionales —los tratados comerciales, la Convención de Compraventa Internacional de Mercaderías, los tratados sobre cables submarinos, sobre aeronaves, sobre contaminación de aguas, sobre pesca en el mar argentino y, entre ellos, los tratados de integración— ocupa el escalón siguiente de la pirámide.

\begin{definition}{Artículo 75, inciso 22, primer párrafo}
``Corresponde al Congreso [\ldots] Los tratados y concordatos tienen jerarquía superior a las leyes.''

\textit{Constitución Nacional Argentina, artículo 75, inciso 22 (primer párrafo).}
\end{definition}

Esta disposición establece que todos los tratados en general —el de integración, el de compraventa internacional de mercaderías, el de transporte, el de contaminación de aguas, el de pesca— tienen jerarquía superior a las leyes internas.

\subsection{El artículo 75 inciso 24 y los órganos supranacionales}

La reforma de 1994 se vincula con el contexto regional de ese mismo momento: culminaba el período de transición del Mercosur y era necesario establecer su régimen de institucionalización permanente, con la expectativa —no concretada— de que sus órganos evolucionaran hacia la supranacionalidad.

\begin{definition}{Artículo 75, inciso 24 de la Constitución Nacional (reforma de 1994)}
Faculta al Congreso a ``aprobar tratados de integración que deleguen competencias y jurisdicción a organizaciones supraestatales en condiciones de reciprocidad e igualdad, y que respeten el orden democrático y los derechos humanos. Las normas dictadas en su consecuencia tienen jerarquía superior a las leyes''.

\textit{Constitución Nacional Argentina, artículo 75, inciso 24.}
\end{definition}

Este inciso otorga al Congreso la facultad de transferir competencias legislativas a órganos de integración supranacional. Las normas dictadas ``en su consecuencia'' —es decir, en consecuencia de la facultad delegada a órganos supranacionales, y no del tratado en sí mismo— tienen jerarquía superior a las leyes. El tratado de integración sigue siendo, en todo caso, el escalón superior.

De este modo queda reconstruida la pirámide normativa argentina, tal como resultó ordenada a partir de la reforma de 1994.

\begin{table}[H]
\centering
\caption{Pirámide normativa argentina posterior a la reforma de 1994}
\begin{tabularx}{\textwidth}{@{}>{\raggedright\arraybackslash}p{0.28\textwidth}X@{}}
\toprule
\textbf{Escalón} & \textbf{Contenido} \\
\midrule
Primer escalón (cúspide) & Constitución Nacional y tratados de derechos humanos con jerarquía constitucional (bloque de constitucionalidad, art. 75 inc. 22). \\
\addlinespace
Segundo escalón & Los demás tratados internacionales, incluidos los tratados de integración (art. 75 inc. 22, primer párrafo). \\
\addlinespace
Tercer escalón (denominado en clase ``cuarto escalón'') & Derecho derivado de los órganos de integración supranacional (art. 75 inc. 24). \\
\addlinespace
Cuarto escalón & Leyes dictadas por el Congreso de la Nación. \\
\bottomrule
\end{tabularx}
\end{table}

% TODO: contenido incompleto o ambiguo en las notas originales — el apunte denomina al derecho derivado ``cuarto escalón'' pese a ubicarlo en el tercer nivel de la enumeración de cuatro escalones; se conserva la denominación tal como figura en la clase.

El derecho derivado de la integración tiene jerarquía superior al Código Civil, al Código Penal, al código procesal y al derecho interno en general —según la materia de que se trate, dado que el Mercosur no legisla, por ejemplo, sobre derecho penal—. Se ubica, así, en un escalón intermedio, justo antes de la ley interna.

\section{El derecho derivado de la integración en la práctica argentina}

\subsection{El caso ``Cafés La Virginia''}

Un segundo antecedente jurisprudencial se refiere específicamente al derecho de la integración, a diferencia de ``Ekmekdjian'', que constituía un caso de derechos humanos.

\begin{example}{El caso ``Cafés La Virginia S.A.'' (Corte Suprema de Justicia de la Nación, 1994)}
La empresa Cafés La Virginia había celebrado un contrato de compraventa de granos de café con una empresa brasileña. Al ingresar la mercadería a la República Argentina, la Aduana le cobró un arancel de importación, aplicable en principio a productos extranjeros para proteger la producción nacional.

Sin embargo, en ese momento se encontraba vigente entre la Argentina y Brasil un acuerdo de alcance parcial, celebrado en el marco de la ALADI (Asociación Latinoamericana de Integración) —no todavía del Mercosur—, que eximía del pago de ese arancel al grano de café.

La Corte Suprema aplicó el mismo principio que en ``Ekmekdjian'': el acuerdo internacional, dictado como consecuencia de un tratado de integración, tenía jerarquía superior a la ley interna —en este caso, a la normativa aduanera que había impuesto el arancel—. La Corte analizó, además, que la norma aduanera había sido dictada por autoridad competente y siguiendo el procedimiento correspondiente —el Congreso había delegado esas facultades en el Poder Ejecutivo y éste en la Aduana—, pero que, aun así, contrariaba una norma de jerarquía superior: el acuerdo celebrado en el marco de la ALADI.
\end{example}

Este antecedente versa sobre derecho de integración propiamente dicho, mientras que ``Ekmekdjian'' se refería a derechos humanos; sin embargo, en ambos casos la Corte sostuvo la misma jerarquía normativa: el acuerdo dictado en el marco de un área de integración prevaleció sobre la norma interna.

\subsection{La jerarquía del derecho del Mercosur según cada Estado}

La jerarquía según la cual el tratado y el derecho de integración se ubican por encima de la ley rige puertas adentro del sistema argentino, pero no necesariamente en los demás Estados parte del Mercosur, porque cada uno cuenta con su propio sistema constitucional de recepción del derecho internacional.

\begin{table}[H]
\centering
\caption{Jerarquía del derecho derivado del Mercosur según el Estado parte}
\begin{tabularx}{\textwidth}{@{}>{\raggedright\arraybackslash}p{0.18\textwidth}>{\raggedright\arraybackslash}p{0.25\textwidth}X@{}}
\toprule
\textbf{País} & \textbf{Sistema} & \textbf{Consecuencia para el derecho derivado del Mercosur} \\
\midrule
Brasil & Dualista & Ni siquiera corresponde hablar de jerarquía normativa: el derecho internacional y el derecho interno se mueven en ámbitos distintos que no se cruzan automáticamente. \\
\addlinespace
Paraguay y Uruguay & Monista, sin cláusula de supremacía & Son monistas, pero no cuentan con una cláusula constitucional como el artículo 75 inciso 24 argentino que otorgue jerarquía superior expresa al derecho derivado. \\
\addlinespace
Argentina & Monista, con cláusula de supremacía (art. 75 inc. 24) & El derecho derivado del Mercosur tiene jerarquía superior a las leyes internas. \\
\bottomrule
\end{tabularx}
\end{table}

Esta supremacía sólo puede afirmarse, en general, para la República Argentina, y no para Paraguay ni para Brasil, ya que en esos países no surge de ninguna norma que el derecho de la integración tenga jerarquía superior al derecho interno. No obstante esta diferencia formal, en la práctica judicial argentina los jueces suelen reconocer supremacía al derecho derivado del Mercosur frente a la ley interna, incluso más allá de lo estrictamente previsto por la norma, lo que constituye una interpretación amplia y, en algunos casos, poco rigurosa.

\section{Principios del derecho de la integración y del derecho comunitario}

\subsection{El modelo de la Unión Europea: monismo con supremacía absoluta}

Existe una tercera posibilidad teórica, distinta de la argentina: un monismo con supremacía absoluta del derecho de la integración, incluso por encima de las constituciones nacionales. Para buena parte de la doctrina, esta sería la situación de la Unión Europea, en la medida en que muchas constituciones nacionales debieron adecuarse para permitir la transferencia de facultades a los órganos comunitarios.

\begin{note}
Esta posición doctrinaria —el derecho comunitario por encima incluso de la Constitución nacional— no es evidente ni pacífica, sino una construcción doctrinaria discutida.
\end{note}

\subsection{El principio de especificidad}

\begin{definition}{Principio de especificidad}
El derecho de la integración solamente puede constar de aquellos temas específicos que son objeto de la integración de que se trate. Si el proceso de integración corresponde a una unión aduanera, sólo puede regularse lo relativo al arancel externo común y a la zona de libre comercio y libre tránsito, sin poder extenderse a otras materias no comprendidas en el tratado constitutivo.
\end{definition}

El principio de especificidad implica que un órgano de integración solamente puede regular aquello que le fue expresamente delegado, sin poder extender su competencia a materias no comprendidas en el tratado constitutivo.

\subsection{Competencias exclusivas: la analogía del mandato y el caso del Brexit}

El derecho derivado de los órganos de integración goza de jerarquía superior a la ley interna porque el propio Estado, en ejercicio de su soberanía, decidió transferir esa competencia y se abstuvo de seguir regulando esa materia. La transferencia de una competencia puede compararse con el otorgamiento de un mandato: quien recibe el poder puede ejercerlo, pero no puede ir en contra de lo que ese poder habilita; y quien lo otorgó conserva, en principio, la potestad de revocarlo.

\begin{example}{El Brexit}
La salida del Reino Unido de la Unión Europea demuestra que la transferencia de competencias no implica una renuncia irreversible a la soberanía. Si la Unión Europea hubiera absorbido la soberanía británica de manera absoluta, el Reino Unido no habría podido retirarse. Pudo hacerlo porque el Estado no renunció a esa competencia, sino que la transfirió. Por ello se habla de que los Estados transfieren competencias en forma exclusiva a los órganos de integración o comunitarios: mientras dura la transferencia, sólo ese órgano puede dictar ese tipo de normas, pero únicamente dentro del ámbito específico delegado, no sobre cualquier materia.
\end{example}

La Unión Europea recibió competencias en materias muy amplias —libre tránsito de mercaderías, personas y servicios, emisión de moneda, medio ambiente, responsabilidad parental, entre otras—, lo que explica que, a simple vista, pareciera situarse ``por arriba de todo''. Sin embargo, en teoría, la Constitución de un Estado miembro nunca queda por debajo del derecho comunitario en un sentido absoluto, porque ese derecho comunitario deriva del poder que le otorgó el Estado soberano a través de su propia Constitución. Un ciudadano argentino que llega a España recibe la aplicación del derecho interno español, y no el derecho de la Unión Europea, porque no es ciudadano europeo: el derecho interno sigue funcionando para quienes no forman parte de la Unión.

A diferencia de la Unión Europea, en la que los órganos actúan por mayoría, en el Mercosur los órganos actúan de manera intergubernamental, lo que evidencia que es el Estado quien sigue tomando la decisión, sin que ésta le sea impuesta.

\section{Aplicación práctica y vigencia del derecho derivado}

\subsection{El principio de vigencia simultánea en el Mercosur}

El carácter intergubernamental del Mercosur se traduce, en la práctica, en el mecanismo de vigencia simultánea del derecho derivado.

\begin{example}{Los primeros protocolos del Mercosur}
En los primeros tiempos del Mercosur, las normas se aprobaban mediante protocolos firmados por los presidentes de los Estados parte. Cada protocolo establecía que entraría a regir en el ámbito del Mercosur a partir de que dos de los Estados hubieran depositado su instrumento de ratificación. El problema era que, con ese sistema, una norma podía estar vigente para dos Estados y no para los otros dos, generando situaciones asimétricas.

Esta asimetría se verificó, por ejemplo, con los tratados de medidas cautelares: la Argentina solicitaba a Brasil, a través del protocolo correspondiente, que no se tomara una medida cautelar sobre determinada empresa por vía de la autoridad central, y Brasil respondía que, en su ordenamiento interno, las medidas cautelares y los exhortos eran de competencia exclusiva del Tribunal Supremo, de modo que esa norma no resultaba vigente para su Estado en el ámbito interno. El resultado era una obligación asimétrica: un Estado quedaba obligado y el otro no.
\end{example}

Para corregir este problema se estableció el principio de vigencia simultánea: la normativa que emana de los órganos del derecho derivado entra en vigor en todos los Estados parte al mismo tiempo, nunca antes. Se necesita consenso, con todos los Estados presentes, para que la norma se dicte; pero la norma existe sin entrar en vigencia hasta que se cumple un procedimiento específico en todos los países.

\begin{table}[H]
\centering
\caption{Procedimiento de vigencia simultánea del derecho derivado del Mercosur}
\begin{tabularx}{\textwidth}{@{}>{\raggedright\arraybackslash}p{0.28\textwidth}X@{}}
\toprule
\textbf{Paso} & \textbf{Descripción} \\
\midrule
1. Incorporación al derecho interno &
La normativa del Mercosur no exige ``ratificación'' del Congreso —propia de un sistema monista puro— sino ``incorporación por parte de los Estados'', término que remite más bien a un sistema dualista. Cada Estado debe cumplir el procedimiento interno que corresponda para incorporar esa norma de integración a su derecho nacional. \\
\addlinespace
2. Notificación a la Secretaría del Mercosur &
Una vez incorporada la norma, el Estado debe notificar ese acto a la Secretaría del Mercosur. \\
\addlinespace
3. Espera de todos los documentos de incorporación &
La Secretaría recibe los documentos de incorporación de cada uno de los Estados parte, pero no comunica nada mientras falte alguno: la norma no entra en vigor para ninguno de los Estados hasta que todos hayan notificado su incorporación. \\
\addlinespace
4. Entrada en vigor simultánea &
Una vez recibido el último documento de incorporación, la Secretaría notifica que la norma entrará en vigor dentro del plazo establecido en cada caso (por ejemplo, dentro de tres o cinco días), y esa entrada en vigor rige para todos los Estados al mismo tiempo. \\
\bottomrule
\end{tabularx}
\end{table}

Este mecanismo demuestra que el Mercosur no es supranacional, sino intergubernamental, y que además cuenta con una etapa de prevención: si un Estado no cumplió el procedimiento de incorporación, la norma adoptada por consenso de todos sigue sin poder aplicarse a ninguno.

La diferencia central respecto del sistema de la Unión Europea radica en que, mientras en el Mercosur la vigencia depende de la incorporación de cada Estado, en la Unión Europea la norma se aplica igual aunque un Estado no la haya votado a favor, porque allí sí se habla, en sentido estricto, de supremacía y de efecto directo del derecho comunitario.

\section{Cuadro Cornell: síntesis del bloque temático}

\begin{longtable}{@{}>{\raggedright\arraybackslash}p{0.28\textwidth}>{\raggedright\arraybackslash}p{0.62\textwidth}@{}}
\caption{Cuadro Cornell --- Jerarquía normativa, derechos humanos y derecho derivado en el Mercosur} \\
\toprule
\textbf{Preguntas / palabras clave} & \textbf{Notas / respuestas} \\
\midrule
\endfirsthead
\toprule
\textbf{Preguntas / palabras clave} & \textbf{Notas / respuestas} \\
\midrule
\endhead
\midrule
\multicolumn{2}{r}{\textit{Continúa en la página siguiente}} \\
\endfoot
\bottomrule
\endlastfoot

\textbf{¿Qué dice el artículo 31 de la Constitución Nacional?} &
Declara que la Constitución, las leyes dictadas por el Congreso y los tratados con potencias extranjeras son la ``ley suprema de la Nación''. Es la norma que, desde 1853, dio origen al debate sobre la jerarquía entre ley y tratado. \\
\addlinespace

\textbf{¿Qué estableció el fallo ``Martín'' (1963)?} &
Que el artículo 31 no distingue entre ley y tratado: ambos son ``ley suprema'' con igual jerarquía. Los conflictos entre ambos se resuelven con los criterios de ley especial y ley posterior. \\
\addlinespace

\textbf{Ley especial vs. ley posterior: ¿cuál deroga y cuál no?} &
La ley especial prevalece sobre la general pero no la deroga (sigue vigente para los demás casos). La ley posterior sí deroga a la anterior, cuando ambas son de igual generalidad. \\
\addlinespace

\textbf{¿Qué derechos estaban en juego en ``Ekmekdjian c. Sofovich''?} &
La libertad de expresión (art. 14 CN), invocada por Sofovich, frente al derecho a réplica, invocado por Ekmekdjian, no enumerado expresamente pero fundado en el art. 33 CN (derechos implícitos) y en el art. 14.1 del Pacto de San José de Costa Rica. \\
\addlinespace

\textbf{¿Por qué la Corte recurrió al artículo 27 de la Convención de Viena?} &
Porque el Pacto de San José de Costa Rica exigía reglamentación interna del derecho a réplica y el Estado argentino no la había dictado. El artículo 27 (pacta sunt servanda) impide invocar el derecho interno —incluida la falta de reglamentación— para incumplir un tratado. \\
\addlinespace

\textbf{¿Qué cambió la reforma constitucional de 1994?} &
Incorporó los tratados de derechos humanos enumerados en el art. 75 inc. 22 con jerarquía constitucional (bloque de constitucionalidad); estableció que todos los demás tratados tienen jerarquía superior a las leyes; y en el art. 75 inc. 24 habilitó delegar competencias legislativas en órganos supranacionales, dando al derecho derivado jerarquía superior a las leyes. \\
\addlinespace

\textbf{¿Cuál es la pirámide normativa argentina resultante?} &
1) Constitución y tratados de derechos humanos con jerarquía constitucional; 2) demás tratados, incluidos los de integración; 3) derecho derivado de órganos de integración supranacional; 4) leyes del Congreso. \\
\addlinespace

\textbf{¿Qué demostró el caso ``Cafés La Virginia''?} &
Que un acuerdo de integración (en el marco de la ALADI) prevalece sobre una norma aduanera interna posterior, aplicando el mismo principio que ``Ekmekdjian'', pero en materia de derecho de la integración propiamente dicho. \\
\addlinespace

\textbf{¿Rige la misma jerarquía en todos los países del Mercosur?} &
No. Depende del sistema de cada Estado: Brasil es dualista; Paraguay y Uruguay son monistas pero sin una cláusula de supremacía como la argentina (art. 75 inc. 24); sólo en la Argentina el derecho derivado tiene jerarquía superior expresa a la ley. \\
\addlinespace

\textbf{¿Qué es el principio de especificidad?} &
Que un órgano de integración sólo puede regular las materias específicas que le fueron delegadas (por ejemplo, arancel externo común y libre tránsito en una unión aduanera), sin poder extender su competencia a otros temas. \\
\addlinespace

\textbf{¿Por qué el Brexit demuestra que la transferencia de competencias no es una pérdida definitiva de soberanía?} &
Porque el Reino Unido pudo retirarse de la Unión Europea: si la transferencia hubiera sido una renuncia absoluta a la soberanía, esa salida no habría sido posible. El Estado transfiere el ejercicio de una competencia, no la titularidad soberana sobre ella. \\
\addlinespace

\textbf{¿Qué es la vigencia simultánea en el Mercosur?} &
El mecanismo por el cual una norma del derecho derivado del Mercosur sólo entra en vigor cuando todos los Estados parte incorporaron esa norma a su derecho interno y lo notificaron a la Secretaría del Mercosur, evitando que rija para unos Estados y para otros no. \\
\addlinespace

\multicolumn{2}{p{0.9\textwidth}}{
\textbf{Resumen:} El sistema jurídico argentino es monista: la Constitución, los tratados y las leyes conviven en un único ordenamiento. Desde el artículo 31 de la Constitución Nacional, la doctrina y la jurisprudencia discutieron durante décadas si un tratado y una ley tenían igual jerarquía —resolviendo sus conflictos con los criterios de ley especial y ley posterior, como en el fallo ``Martín''— hasta que el caso ``Ekmekdjian c. Sofovich'' (1992) modificó ese criterio: invocando el artículo 27 de la Convención de Viena sobre el Derecho de los Tratados, la Corte Suprema reconoció que un tratado internacional no puede quedar subordinado al silencio o la inacción del derecho interno. Ese antecedente inspiró la reforma constitucional de 1994, que en el artículo 75 incisos 22 y 24 ordenó la pirámide: la Constitución y los tratados de derechos humanos con jerarquía constitucional en la cúspide; luego los demás tratados, incluidos los de integración; después el derecho derivado de los órganos de integración (como las normas del Mercosur); y, por último, las leyes del Congreso. El caso ``Cafés La Virginia'' confirmó este orden para el derecho de la integración en sentido estricto. Sin embargo, esa jerarquía sólo rige puertas adentro de la Argentina: cada Estado parte del Mercosur decide, con su propio sistema —dualista en Brasil, monista sin esta cláusula en Paraguay y Uruguay—, qué lugar le da al derecho derivado, lo que explica por qué el Mercosur funciona de manera intergubernamental y no supranacional como la Unión Europea, y por qué se adoptó el mecanismo de vigencia simultánea: ninguna norma del derecho derivado entra en vigor hasta que todos los Estados la incorporaron a su derecho interno y lo notificaron a la Secretaría. Esta arquitectura —anclada en los principios de especificidad y de atribución de competencias exclusivas, que explican por qué un Estado no pierde su soberanía al delegar una facultad, tal como lo demostró el Brexit— resulta indispensable para argumentar, litigar o asesorar sobre la aplicación de normas internacionales, comunitarias o de integración en la Argentina.
} \\

\end{longtable}

\end{document}
